\chapter{\eh{desktop-computer} Capa 7 --- Aplicación (Application Layer)}

\begin{concepto}
La capa de aplicación es la que \textit{ves y usas} directamente. Aquí viven
los protocolos que usan tus apps: HTTP (web), DNS (nombres de dominio),
DHCP (IPs automáticas), SMTP (correo electrónico), FTP (archivos).
\end{concepto}

\section{\eh{satellite-antenna} DHCP --- Dynamic Host Configuration Protocol (RFC 2131)}

\begin{concepto}
\textbf{DHCP} asigna automáticamente configuración de red a los dispositivos
cuando se conectan: dirección IP, máscara de subred, gateway por defecto y
servidores DNS.

\emoji{light-bulb} \textbf{Analogía:} Llegar a un hotel (red) y que la recepción
(servidor DHCP) te asigne automáticamente: número de habitación (IP),
llave maestra del piso (máscara), la dirección del hotel para recibir correo
(gateway) y el directorio telefónico del hotel (DNS). Todo automático, sin
que tú tengas que pedirlo o configurarlo.
\end{concepto}

\subsection{Proceso DORA}

\begin{center}
\begin{tikzpicture}[
    msg/.style={-{Stealth}, thick},
    server/.style={font=\bfseries\small, text=textlight}
]
\node[server] at (0,0) {Cliente};
\node[server] at (8,0) {Servidor DHCP};
\draw[thick, draw=gray!50!white] (0,-0.3) -- (0,-7.5);
\draw[thick, draw=gray!50!white] (8,-0.3) -- (8,-7.5);

\draw[msg, cyan!70!white] (0,-1.0) -- (8,-1.8)
    node[midway, above, sloped, font=\small\bfseries, text=cyan!80!white] {DISCOVER};
\node[font=\tiny, right, text=gray!70!white] at (8,-1.8)
    {Broadcast: ``¿Hay un servidor DHCP?''};

\draw[msg, green!60!white] (8,-2.3) -- (0,-3.1)
    node[midway, above, sloped, font=\small\bfseries, text=green!70!white] {OFFER};
\node[font=\tiny, left, text=gray!70!white] at (0,-3.1)
    {``Te ofrezco 192.168.1.100 por 24h''};

\draw[msg, orange!80!white] (0,-3.6) -- (8,-4.4)
    node[midway, above, sloped, font=\small\bfseries, text=orange!80!white] {REQUEST};
\node[font=\tiny, right, text=gray!70!white] at (8,-4.4)
    {Broadcast: ``Acepto la oferta de 192.168.1.100''};

\draw[msg, red!70!white] (8,-4.9) -- (0,-5.7)
    node[midway, above, sloped, font=\small\bfseries, text=red!80!white] {ACK};
\node[font=\tiny, left, text=gray!70!white] at (0,-5.7)
    {``Confirmado, es tuya hasta las 14:00 de mañana''};

\node[font=\small\itshape, text=green!60!white] at (4,-6.8)
    {\emoji{check-mark-button} Cliente ya tiene IP, máscara, gateway y DNS};
\end{tikzpicture}
\end{center}

\textbf{Información que provee DHCP:}
\begin{itemize}
    \item IP address y máscara de subred
    \item Default gateway
    \item Servidores DNS primario y secundario
    \item Lease time (tiempo de vigencia de la IP)
    \item NTP servers, TFTP server (opcional)
\end{itemize}

\begin{cuidado}
\textbf{Rogue DHCP Server:} si un atacante conecta su propio servidor DHCP,
puede asignarse a sí mismo como gateway o DNS → intercepta todo el tráfico.

Defensa: \textbf{DHCP Snooping} en el switch --- solo permite respuestas DHCP
de puertos confiables (donde está el servidor legítimo).
\end{cuidado}

\section{\eh{globe-with-meridians} DNS --- Domain Name System (RFC 1034/1035)}

\begin{concepto}
\textbf{DNS} traduce nombres de dominio (\texttt{google.com}) a direcciones IP
(\texttt{142.250.80.46}). Sin DNS, tendrías que memorizar IPs para cada sitio.

\emoji{light-bulb} \textbf{Analogía:} DNS es la agenda de contactos de Internet.
No necesitas memorizar ``142.250.80.46'' --- solo escribes ``google.com'' y
DNS busca el número IP por ti, igual que buscar ``mamá'' en tus contactos
en vez de marcar el número completo.
\end{concepto}

\subsection{Jerarquía DNS}

\begin{lstlisting}
  . (root) ── gestionado por IANA, 13 root servers en todo el mundo
  │
  ├── .com  .org  .mx  .net  .io  ...  (TLDs - Top Level Domains)
  │         │
  │      google.com  (Second-level domain, registrado por Google)
  │         │
  │      www.google.com   maps.google.com   api.google.com
  │      (Subdominios, configurados por Google)
  │
  └── .edu  .gob  ... (TLDs especiales)
\end{lstlisting}

\subsection{Tipos de registros DNS}

\begin{table}[h!]
\centering
\renewcommand{\arraystretch}{1.3}
\begin{tabular}{lll}
\toprule
\textbf{Tipo} & \textbf{¿Qué guarda?} & \textbf{Ejemplo} \\
\midrule
A     & IPv4 del dominio           & \texttt{google.com → 142.250.80.46} \\
AAAA  & IPv6 del dominio           & \texttt{google.com → 2607:f8b0::...} \\
CNAME & Alias a otro dominio       & \texttt{www → google.com} \\
MX    & Servidor de correo         & \texttt{google.com → aspmx.l.google.com} \\
TXT   & Texto libre                & SPF, DKIM, verificación de dominio \\
NS    & Nameservers del dominio    & \texttt{google.com → ns1.google.com} \\
PTR   & IP → nombre (reverse DNS) & \texttt{46.80.250.142.in-addr.arpa → google.com} \\
SOA   & Registro de autoridad      & TTL, email del admin, número de serie \\
SRV   & Servicio y puerto          & \texttt{\_http.\_tcp.example.com} \\
\bottomrule
\end{tabular}
\caption{Tipos de registros DNS}
\end{table}

\subsection{Proceso de resolución DNS}

\begin{lstlisting}
  Browser pide: www.unam.mx

  1. Caché del OS → ¿ya lo tengo? Si sí, listo.
  2. Resolver local (tu router / ISP)
  3. Root server → ``No sé, pregunta al TLD de .mx''
  4. TLD server (.mx) → ``No sé, pregunta al NS de unam.mx''
  5. Authoritative NS (ns.unam.mx) → ``La IP es 132.248.204.100''
  6. Resolver lo cachea y te responde
  7. Tú te conectas a 132.248.204.100

  Resolución iterativa: el resolver hace todas las preguntas
  Resolución recursiva: el root server haría todas las preguntas (raro)
\end{lstlisting}

\section{\eh{globe-showing-americas} HTTP --- HyperText Transfer Protocol}

\subsection{HTTP/1.1 --- El original (RFC 2616, 1999)}

\begin{lstlisting}
  ─── Request ─────────────────────────────────────────
  GET /index.html HTTP/1.1
  Host: example.com
  User-Agent: Mozilla/5.0 (Windows NT 10.0)
  Accept: text/html, application/json
  Connection: keep-alive

  ─── Response ─────────────────────────────────────────
  HTTP/1.1 200 OK
  Content-Type: text/html; charset=utf-8
  Content-Length: 2048
  Cache-Control: max-age=3600

  <!DOCTYPE html>...
\end{lstlisting}

\textbf{Problema principal de HTTP/1.1:}
\begin{itemize}
    \item \textbf{Head-of-line blocking:} con pipelining, si la primera petición
          tarda, todas las siguientes esperan en fila.
    \item Una página web con 100 recursos (CSS, JS, imágenes) requiere múltiples
          conexiones TCP (hasta 6 por dominio en browsers modernos).
    \item Cada conexión TCP tiene overhead: handshake + slow start.
\end{itemize}

\subsection{HTTP/2 (RFC 7540, 2015)}

\begin{concepto}
Google desarrolló SPDY (2009) para resolver las limitaciones de HTTP/1.1.
SPDY se convirtió en la base de \textbf{HTTP/2}.
\end{concepto}

Mejoras clave:
\begin{itemize}
    \item \textbf{Multiplexing:} múltiples requests en \textit{una sola} conexión TCP,
          en paralelo real. Sin cola, sin esperas.
    \item \textbf{Header compression (HPACK):} los headers se comprimen y los
          repetidos se envían como diferencias. HTTP/1.1 enviaba todos los headers
          en texto plano en cada request --- enorme overhead.
    \item \textbf{Binary framing:} protocolo binario, no texto. Más eficiente.
    \item \textbf{Server push:} el servidor puede enviar recursos antes de que
          el cliente los pida (enviar CSS y JS mientras el HTML se procesa).
    \item \textbf{Stream prioritization:} el cliente puede indicar qué recursos
          son más urgentes.
\end{itemize}

\subsection{Códigos de estado HTTP}

\begin{table}[h!]
\centering
\renewcommand{\arraystretch}{1.3}
\begin{tabular}{lll}
\toprule
\textbf{Rango} & \textbf{Categoría} & \textbf{Ejemplos importantes} \\
\midrule
1xx & Informacional  & 100 Continue, 101 Switching Protocols \\
2xx & Éxito          & 200 OK, 201 Created, 204 No Content \\
3xx & Redirección    & 301 Moved Permanently, 304 Not Modified \\
4xx & Error cliente  & 400 Bad Request, 401 Unauthorized, 403 Forbidden, 404 Not Found \\
5xx & Error servidor & 500 Internal Server Error, 502 Bad Gateway, 503 Service Unavailable \\
\bottomrule
\end{tabular}
\caption{Códigos de estado HTTP}
\end{table}

\section{\eh{envelope} Correo Electrónico --- SMTP, POP3, IMAP}

\begin{table}[h!]
\centering
\renewcommand{\arraystretch}{1.3}
\begin{tabular}{llll}
\toprule
\textbf{Protocolo} & \textbf{Puerto} & \textbf{Uso} & \textbf{Dirección} \\
\midrule
SMTP     & 25, 587 (TLS) & Enviar correo             & Cliente → Servidor, Servidor → Servidor \\
POP3     & 110, 995 (TLS) & Recibir y descargar       & Descarga y borra del servidor \\
IMAP     & 143, 993 (TLS) & Recibir sincronizado      & Sincroniza entre dispositivos \\
\bottomrule
\end{tabular}
\caption{Protocolos de correo electrónico}
\end{table}

\section{\eh{file-folder} FTP y sus alternativas}

\begin{table}[h!]
\centering
\renewcommand{\arraystretch}{1.3}
\begin{tabular}{llll}
\toprule
\textbf{Protocolo} & \textbf{Puerto} & \textbf{Cifrado} & \textbf{Estado} \\
\midrule
FTP   & 21 (control), 20 (datos) & \emoji{cross-mark} No   & Evitar, inseguro \\
FTPS  & 990                      & \emoji{check-mark-button} TLS & Obsolescente \\
SFTP  & 22 (sobre SSH)           & \emoji{check-mark-button} SSH & Recomendado \\
SCP   & 22 (sobre SSH)           & \emoji{check-mark-button} SSH & Recomendado \\
\bottomrule
\end{tabular}
\caption{Protocolos de transferencia de archivos}
\end{table}

\begin{cuidado}
FTP envía usuario y contraseña en texto claro. \textbf{Nunca usar FTP}
para transferencias reales. SFTP usa el canal SSH (puerto 22) y es
completamente diferente a FTPS (que es FTP + TLS).
\end{cuidado}
